\section{A. The Deuring element and its local comparisons}
\label{sec:deuring}

\paragraph{Conventions.}
We write $k=\overline{\mathbb F}_p$, $W=W(k)$, and $\sigma$ for Witt
Frobenius. An integral $F$-crystal has a finite free $W$-module and an
injective $\sigma$-linear operator $F$. Its Tate determinant is
$T=W t$, $F(t)=pt$. Thus $T$ is the crystalline realization of the
cohomological object $\mathbb Z_p(-1)$. We use the covariant convention
\[
 D(E)=H^1_{\rm cris}(E/W),\qquad D(f)=(f^\vee)^*.
\]
Here the canonical principal polarizations identify an elliptic curve
with its dual, so $f^\vee:E'\to E$ for $f:E\to E'$. This convention makes
$D(E)$ a left $\operatorname{End}(E)$-module. It does not change its
$\sigma$-linear crystalline Frobenius. The alternating cup product followed
by the trace is denoted $b_E:\bigwedge^2D(E)\xrightarrow{\sim}T$.
It satisfies
\begin{equation}\label{eq:A-degree}
 b_{E'}\circ\bigwedge^2D(f)=(\deg f)b_E.
\end{equation}
Indeed $D(f)=(f^\vee)^*$ and pullback on $H^2$ multiplies the trace by
$\deg(f^\vee)=\deg f$. This proves the sign and covariance used below.
All gerbes and torsors over rings in this section use the fppf topology;
this matters at $2$. A right projective $O$-module $P$ is acted on from the
left by $O'=\operatorname{End}_O(P)$.

\subsection{The domain of the question, and an actual counterexample}

\paragraph{A.1a (Refuted: the unrestricted $\mathbb F_{p^2}$ formulation).}
The phrase in 15:54, ``un F-cristal $M$ sur $\mathbb F_{p^2}$ (ou sur
$\overline{\mathbb F}_p$) de rang $2$, muni d'une opération de $O$'',
does not by itself ensure that the displayed set of determinant
isomorphisms is nonempty. This remains false after adding slope $1/2$.

\begin{proof}
Take $p=3$, $W_2=W(\mathbb F_9)=\mathbb Z_3[a]$, $a^2=-1$, so
$\sigma(a)=-a$. Put
\[
 J=\begin{pmatrix}0&-3\\1&0\end{pmatrix},\quad
 I=\begin{pmatrix}a&0\\0&-a\end{pmatrix},\quad
 M=W_2^2,\quad F=4J\sigma.
\]
The rational quaternion algebra $B=(-1,-3)_{\mathbb Q}$ has generators
$i,j$ represented by $I,J$. We have $I^2=-1$, $J^2=-3$ and $IJ=-JI$.
The order
\[
 O=\mathbb Z\langle 1,i,(1+j)/2,(i+ij)/2\rangle
\]
is closed under multiplication and its reduced-trace Gram determinant is
$-9$. The Hilbert symbol is $-1$ at $3$ and $\infty$, and $+1$ elsewhere
(at $2$ use reciprocity). Thus $B=B_3$ and reduced discriminant $3$
makes $O$ maximal. These matrix, multiplication and discriminant
calculations are reproduced by \texttt{check\_quaternion\_crystal.py}.
Both $I$ and $J$ commute with the semilinear $F$; integrality of the
half-integral generators holds because $2$ is a unit in $W_2$.
Now $F^2=-48$ and $v_3(48)=1$, so both slopes are $1/2$.
Nevertheless $\bigwedge^2F=48\sigma$. An isomorphism of its determinant
with $(W_2,3\sigma)$ would be multiplication by a unit $c$ satisfying
$\sigma(c)=16c$. Applying $\sigma$ twice gives $c=256c$, impossible in
the characteristic-zero domain $W_2$. In particular there is no rule on
every pair in the literal finite-field domain of A.1.
\end{proof}

\paragraph{A.1b (Proved: the exact determinant qualification).}
For a rank-two $F$-crystal over $\mathbb F_{p^2}$ with a full maximal
$O$-action and slope $1/2$, one has $F^2=pc$ for $c\in\mathbb Z_p^*$.
Its determinant is Tate if and only if $c^2=1$, equivalently $F^2=p$ or
$F^2=-p$. Over $k$ the determinant is always Tate.

\begin{proof}
Over $W_2[1/p]$, $O\otimes\mathbb Q_p$ splits and spans the full
two-by-two matrix algebra. Its commutant consists of scalar matrices.
Consequently $F^2=b$ is scalar; commutation with $F$ gives
$\sigma(b)=b$, and the slope gives $b=pc$ as asserted. The square of
determinant Frobenius is $p^2c^2$, whereas that of $T$ is $p^2$.
This proves necessity. For sufficiency write the determinant operator
as $pd\sigma$; then $d\sigma(d)=1$. Integral Hilbert 90 for the
unramified quadratic extension gives $d=\sigma(e)/e$ with
$e\in W_2^*$: ordinary Hilbert 90 gives an element of its fraction
field and multiplication by a power of $p$ makes it a unit.
Over $W(k)$ the equation $\sigma(e)/e=d$ for a unit is solved successively
modulo $p^r$: the first equation is $x^{p-1}=\bar d$, and each subsequent
correction is an additive equation $x^p-x=t$, soluble in $k$.
Completeness gives a unit solution. No division by $2$ is used.
\end{proof}

\subsection{A precise geometric rule over the algebraic closure}

\paragraph{A.1c (Proved: local classification needed by the rule).}
Every rank-two slope-$1/2$ integral $F$-crystal over $k$ is isomorphic to
\[
 M_* =We\oplus Wf,\qquad Fe=f,\quad Ff=-pe.
\]
Its endomorphism ring as an $F$-crystal is the maximal order in the
quaternion division algebra over $\mathbb Q_p$. Two full actions of a
fixed $O$ on such crystals are integrally conjugate precisely when their
residue orientations agree. An $O$-linear isomorphism between them, if
it exists, is unique up to $\mathbb Z_p^*$.

\begin{proof}
The determinant has valuation one. Thus the elementary divisors of $F$
are $1,p$ and $V=pF^{-1}$ preserves the lattice. The rational
Dieudonné--Manin classification gives the isocrystal $M_*[1/p]$
(\citealp[Theorem 5]{Yu2026}). Here is the integral step, often suppressed.
Choose an embedding of the given lattice in $M_*$ with minimal index;
one exists by scaling an isocrystal isomorphism by a power of $p$.
Its image modulo $p$ is nonzero, since otherwise dividing the embedding
by $p$ would reduce the index. If it is not onto, that image is the unique line
stable under the nilpotent rank-one operators $F,V$, namely $kf$.
The inverse image $pWe\oplus Wf$ is itself isomorphic to $M_*$ by
$e\mapsto f$, $f\mapsto-pe$. The index in this lattice is smaller,
a contradiction. This also proves the assertion for all primes,
including $2$.

Writing a commuting matrix out shows that the endomorphism ring has
the presentation $W(\mathbb F_{p^2})\oplus W(\mathbb F_{p^2})\Pi$,
$\Pi^2=-p$, $\Pi a=\sigma(a)\Pi$. It is the valuation ring of the
division algebra and hence maximal. By Skolem--Noether two embeddings
are conjugate by a division-algebra unit in the rational sense. Writing
that conjugator as $u\Pi^r$ shows that its residue action is Frobenius
to the power $r$; the orientations agree exactly when $r$ is even.
Multiplication by the central $p^{-r/2}$ then makes the conjugator an
integral unit. Finally, the centralizer of the full quaternion action
is scalar, and commutation with $F$ cuts the integral scalar units down
to $W(k)^{\sigma,*}=\mathbb Z_p^*$.
\end{proof}

\paragraph{A.1d (Proved: geometric determination of the Deuring rule).}
For a pair $(M,O)$ over $k$, define an orientation
$\zeta_M:O\to k$ by the action of $a$ on $M/FM$. Choose a supersingular elliptic curve $E/k$ and
an oriented isomorphism $\eta:O\xrightarrow{\sim}\operatorname{End}(E)$,
and choose an $O$-linear crystalline isomorphism
$b:D(E)\xrightarrow{\sim}M$. Then
\begin{equation}\label{eq:A-Deuring}
 \boxed{\quad u_D(M,O)=
 [\,b_E\circ\bigwedge^2 b^{-1}\,]
 \ \in\operatorname{Isom}_{F}(\bigwedge^2M,T)/\mathbb Z_p^{*2}.
 \quad}
\end{equation}
The right side exists and is independent of all these choices. It is
the rule induced by the geometric realization in 15:50(a). It satisfies
(i) and (ii) of 15:54, including field automorphisms. All such rules form
a simply transitive set under $\mathbb Z_p^*/\mathbb Z_p^{*2}$: two
elements if $p$ is odd, and four if $p=2$. Lead A1 is correct on this
qualified domain; it is not a repair of A.1a without changing its domain.

\begin{proof}
For a supersingular elliptic curve, the image of crystalline Frobenius
modulo $p$ is the Hodge line $H^0(E,\Omega^1)$: Frobenius kills
that line and its image in $H^1(E,\mathcal O_E)$ is the zero
Hasse--Witt map, while its rank is one. Therefore
$D(E)/FD(E)=H^1(E,\mathcal O_E)=\operatorname{Lie}(E^\vee)$.
The action $(a^\vee)^*$ on this quotient is the differential of
$(a^\vee)^\vee=a$. Thus $\zeta_{D(E)}$ is exactly Voight's
tangent-space orientation. This argument fixes the convention directly,
without identifying crystalline Frobenius with an unspecified
Dieudonné convention.
The oriented Deuring theorem (\citealp[Proposition 42.4.8]{Voight2021})
supplies $(E,\eta)$, uniquely up to an isomorphism whose two possible
choices differ by $[-1]$. A.1c supplies $b$. Replacing $b$ by $cb$,
$c\in\mathbb Z_p^*$, multiplies the displayed determinant isomorphism
by $c^{-2}$. Changing the marked curve changes it by the determinant
of a curve isomorphism, which is its degree, namely one. This proves
independence, and also proves invariance under isomorphisms of pairs.

For a right projective $O$-module $P$, the covariant Deuring equivalence
identifies $P\otimes_OE$ with an elliptic curve $E'$ and crystalline
realization identifies its crystal with $P\otimes_OD(E)$. The determinant
comparison is the norm-line comparison. More explicitly, identify
$P$ with $\operatorname{Hom}(E,E')$. For a nonzero rational generator
$t$ and $a\in\operatorname{End}^0(E)$,
$\deg(ta)=\deg(t)\operatorname{Nrd}(a)$; hence degree defines the same
transition functions as the reduced norm on $P$. This norm line maps
isomorphically to $\mathbb Z$, not to a proper ideal: at every
$\ell\ne p$ the integral Tate homomorphism theorem, over a common
finite field with scalar Frobenius, supplies an invertible map of
the two Tate lattices; at $p$ the integral crystalline version and
A.1c do the same. These integral homomorphism theorems are recorded
in \citet[\S1.1--1.2, pp.\ 524--525]{Waterhouse1969}.
Approximating each such local map by a global element of
$\operatorname{Hom}(E,E')$ yields an isogeny of degree prime to that
particular $\ell$. Thus the degrees generate the unit ideal in
$\mathbb Z$. Positivity makes the generator $+1$.
By \eqref{eq:A-degree} this is exactly the norm trivialization of
A.6a, so the diagram in 15:54 commutes without a hidden unit factor.
This proves (ii), including changes of the right-module presentation.
Cup product, trace, duality and the positive degree are invariant under
every automorphism of $k$; hence (i) includes Frobenius.

For any other rule $u$, the quotient $u/u_D$ is a unit squareclass.
Every pair is realized as above and the Deuring correspondence connects
any two marked curves by a rank-one invertible bimodule. Condition (ii)
therefore makes this quotient constant. Conversely a constant unit
squareclass multiplies every value and preserves both conditions.
For odd $p$, reduction identifies unit squareclasses with those in
$\mathbb F_p^*$; for $p=2$ an odd unit is a square precisely when it
is $1$ modulo $8$. This gives the asserted cardinalities.
\end{proof}

\paragraph{A.1e (Proved: an evaluation prescription and its scope).}
Formula \eqref{eq:A-Deuring} determines a choice-independent output.
One way to evaluate it is to enumerate supersingular
Weierstrass models over finite subfields, enumerate endomorphisms with
the degrees of a chosen integral basis of $O$, test the multiplication
table and the orientation, and solve the $O,F$-intertwining equations
for $b$. The search terminates by A.1c and the oriented Deuring theorem.
Choose a symplectic crystalline basis $(x,y)$ of $D(E)$ with
$b_E(x\wedge y)=t$; if $b$ has matrix $B$ in a basis of $M$, the
coefficient in that basis is $(\det B)^{-1}$, taken modulo squares.
The pairing can be evaluated by the de Rham trace on any smooth lift
and the residue pairing at the origin. Different lifts and symplectic
bases give the same expression by the proof of A.1d. More explicitly,
first solve the $O$-intertwining equations for an integral module
isomorphism $C:D(E)\to M$. The matching residue orientation and A.1c
ensure that such a $C$ exists. The two semilinear operators differ by
a scalar unit: $F_M C=aC F_{D(E)}$. Solve
$\sigma(r)/r=a^{-1}$ in $W(k)^*$ by the successive equations in
A.1b; then $b=Cr$ is the required $F$-intertwiner. All its choices
differ by $\mathbb Z_p^*$ and hence disappear in the specified
square quotient. This is a prescription on the input crystal and
oriented order, with the geometric curve recovered by the preceding
terminating search; it is not a choice of an extra polarization on
the input. We have not implemented this all-$p$ crystalline evaluation algorithm.
The executable certificates in this package establish the explicit
counterexamples and finite counts, not numerical values of this
geometric formula for arbitrary input orders. No unrun computation
is used as a verification event.

\paragraph{A.1h (Proved: descent of the rule on its exact finite-field domain).}
Let $(M,O)$ be as in A.1b over $W_2=W(\mathbb F_{p^2})$, and suppose
$F^2=p$ or $F^2=-p$. Then the rule of A.1d descends uniquely to
\[
 u_{D,2}(M,O)\in
 \operatorname{Isom}_{F,W_2}(\bigwedge^2M,T)/\mathbb Z_p^{*2}.
\]
It is independent of the embedding $\mathbb F_{p^2}\hookrightarrow k$
and satisfies the same isomorphism and Morita compatibilities.
This gives the rule over the largest finite-field domain on which
the target torsor in 15:54 is nonempty, including $p=2,3$.

\begin{proof}
A.1b makes the source of the base-change map
\[
 \operatorname{Isom}_{F,W_2}(\bigwedge^2M,T)
 \longrightarrow
 \operatorname{Isom}_{F,W}(\bigwedge^2(M\otimes_{W_2}W),T)
\]
nonempty. Both sides are torsors under the same group $\mathbb Z_p^*$:
an automorphism of the Tate $F$-crystal is a unit fixed by $\sigma$,
over either coefficient ring. The map is equivariant for the identity
of this group, hence bijective. The same holds after quotienting by
squares. Transporting A.1d back through this bijection defines the
displayed class. Two embeddings of the finite field differ by
Frobenius; the field-automorphism invariance proved in A.1d identifies
the resulting classes. Isomorphism and Morita compatibility follow
after this injective base change, so hold before it.

If one studies only the finite-field category and does not require
compatibility with extension to $k$, the two signs of $F^2$ are
separate components. Morita transport preserves the scalar $F^2$.
An $F$-intertwiner between opposite signs would satisfy $pb=-pb$
and is zero, since the coefficient ring has characteristic zero.
Thus constants on the two sign strata can be chosen independently.
The geometric base-change condition is what selects the common
normalization above; it must not be omitted from an assertion that
all finite-field rules form just one unit-squareclass torsor.
\end{proof}

\subsection{The alleged equivalence with quadratic-order triples}

\paragraph{A.1f (Refuted; corrected descent statement).}
The sentence in 15:56 identifying an embedding of the \emph{maximal}
order $R$ of $\mathbb Q(\sqrt{-p})$ with every descent having
$\pi^2+p=0$ is false. The order forced by the equation is
$R_0=\mathbb Z[\sqrt{-p}]$. For $p\equiv3\pmod4$ it extends to $R$
exactly when the resulting $\mathbb F_p$-curve has all its geometric
$2$-torsion rational. For $p=2$ or $p\equiv1\pmod4$, $R_0=R$.

\begin{proof}
Take $E:y^2=x^3+x$ over $\mathbb F_7$. Its discriminant is $6$ modulo
$7$ and its affine points are
\[
 (0,0),(1,3),(1,4),(3,3),(3,4),(5,2),(5,5).
\]
Thus $\#E(\mathbb F_7)=8$, its trace is zero, it is supersingular and
$\pi^2=-7$. Its $2$-torsion has $y=0$ and
$x(x^2+1)=0$; $-1$ is not a square in $\mathbb F_7$, so only one
nonzero $2$-torsion point is rational. If $(1+\pi)/2$ were an
endomorphism, then $1+\pi$ would kill all of $E[2]$, contrary to this
calculation. Conversely, for odd $p$, killing $E[2]$ is precisely the
condition for a homomorphism to factor through $[2]$, since $[2]$ is
the quotient by that finite étale subgroup. This proves both directions.
The exhaustive program \texttt{check\_supersingular.py} checks the
listed points and the small-prime boundary cases. For $p=3$ the same
distinction occurs between $y^2=x^3-x$ and $y^2=x^3+x$.
An arbitrary crystalline descent need not even put its Frobenius in
the global order $O$, only in $O\otimes\mathbb Z_p$; the global-integrality
condition must also be specified. Thus the corrected statement concerns
geometric descents with Frobenius in $O$, not every local descent datum.
\end{proof}

\subsection{The determinant normalization}

\paragraph{A.6a (Proved: the positive, or $\Phi$-normalized, determinant).}
Let $A$ be a Dedekind domain, $K$ its fraction field, $O$ a maximal
order in a central simple $K$-algebra $B$, and $\Phi=\operatorname{Nrd}(B^*)$.
If
\[
 \operatorname{Pic}(A)=0,\qquad \Phi\cap A^*=\{1\},
 \qquad \Phi A^*=K^*,
\]
the norm line of every rank-one right projective $O$-module has a
unique generator in its distinguished $\Phi$-coset. Sending that
generator to $1$ defines natural, multiplicative trivializations of
norm lines. These give the transitive determinant identifications of
15:21--22. For $A=\mathbb Z$, $B=B_p$, the distinguished generator
is the positive generator, and the construction exists for every $p$.

\begin{proof}
Choose an identification $P_K\simeq B$ and view $\operatorname{Nrd}(P)$
as an invertible fractional $A$-ideal $I$. Its distinguished coset in
$I_K^*$ is the image of the norm torsor of bases of $P_K$.
Changing a base multiplies both the line coordinate and this coset by
an element of $\Phi$. Since $I$ is principal, write $I=aA$.
The two group conditions give a unique $u\in A^*$ such that $au$
belongs to the distinguished coset. This generator is independent of
the coordinate choice. Isomorphisms of modules preserve the coset,
so preserve the generator. Tensor product multiplies both norm lines
and distinguished cosets; uniqueness therefore proves multiplicativity,
including the associativity constraint. Applied to
$\det(P\otimes_OX)=\operatorname{Nrd}(P)\otimes\det X$ it gives
the required transitive system.

For $B_p$, the reduced norm is positive at the real ramified place,
and every positive rational is a reduced norm by Hasse--Schilling
(\citealp[Main Theorem 14.7.4]{Voight2021}). Thus
$\Phi=\mathbb Q_{>0}^*$, and the three conditions follow.
For an elliptic realization the positive norm is its degree, so
\eqref{eq:A-degree} identifies the common determinant object with the
Tate object through the principal polarizations. Its crystalline
comparison is exactly A.1d. This proves the precise content of Lead A4.
\end{proof}

\paragraph{A.6b (Refuted: uniqueness without the normalization).}
The three arithmetic hypotheses do not make an arbitrary transitive,
isomorphism-natural system of determinant identifications unique unless
the distinguished norm coset is required. If $t_{XY}$ is one such
system and $s_X\in\{\pm1\}$ is constant on isomorphism classes, then
$t'_{XY}=s_Yt_{XY}s_X^{-1}$ is another. Over $\overline{\mathbb F}_{11}$,
the curves with $j=0$ and $j=1728$ have automorphism groups of orders
$6$ and $4$ respectively. Taking $s_X=-1$ in the first class and $+1$
in the second gives a distinct system, natural even under equivalences
that preserve endomorphism rings. The positive construction in A.6a
removes this ambiguity. There is also still a scalar automorphism of
the common determinant line; a canonical object is not thereby a
uniquely trivialized line.

\subsection{The adelic torsor, its band, and the real sign}

\paragraph{A.2a (Proved: the base ring and the group calculation).}
The ring denoted $\mathbb A$ on 15:46 is explicitly
\[
 A_{\rm int}=\mathbb R\times\prod_\ell\mathbb Z_\ell,
 \qquad A'=\prod_{\ell\ne p}\mathbb Z_\ell.
\]
It is not the usual ring of rational adèles. With fppf cohomology,
\[
 H^1(A_{\rm int},\mu_2)
 =\mathbb R^*/\mathbb R^{*2}\times
   \prod_\ell\mathbb Z_\ell^*/\mathbb Z_\ell^{*2}.
\]
If $U=A'^*/A'^{*2}$, $V=\mathbb Z_p^*/\mathbb Z_p^{*2}$ and
$W_\infty=\mathbb R^*/\mathbb R^{*2}$, then
\[
 \Gamma'=U\times(V\times\mathbb Z/2)\times W_\infty,
 \qquad f=([p],(1,1),[p]),
\]
where $1$ in $V$ is its neutral element (the manuscript uses additive
$0$). The homomorphism
\[
 (u,v,e,w)\longmapsto(u[p]^{-e},v,w[p]^{-e})
\]
has kernel $\langle f\rangle$ and identifies the quotient with
$\Gamma'_0=U\times V\times W_\infty$. In particular
$\Gamma'=\Gamma'_0\times\langle f\rangle$.

\begin{proof}
Every invertible module over a product of local rings is free. For
example, present it by a finite idempotent matrix, choose a generator
coordinatewise, and use the finitely many possible nonzero minors to
partition the index set by central idempotents; these generators then
glue to a basis over the product. The fppf Kummer sequence and this
vanishing of the Picard group give the displayed $H^1$.
At a split prime an invertible bimodule is a scalar ideal and the
determinant rigidification contributes its norm modulo squares.
At $p$, every two-sided ideal is a power of the unique prime
$\mathfrak m$, with $\mathfrak m^2=pO_p$; its norm line is
$p\mathbb Z_p$. The pair $(\mathfrak m,p)$ has order two after
rigidification, while the remaining determinant choices give $V$.
Away from $p$, $\mathfrak m$ is the trivial ideal with determinant
rigidification $p$, giving $[p]$. At infinity $p$ is positive and hence
a square. This proves the claimed local image of $f$. The kernel and
splitting follow immediately from its parity coordinate.
This independently checks R2, with the additive/multiplicative
notation made explicit.
\end{proof}

\paragraph{A.2b (Proved under the geometric comparison interpretation).}
For a positively determinant-rigidified Deuring category $C$ of type
$B_p$, take isomorphism classes of comparison equivalences with the
prime-to-$p$ Tate category, the crystalline category, and the real
quaternionic category. Their product $\mathcal T'(C)$ is a right
$\Gamma'$-torsor. Global precomposition uses the unrigidified
$\operatorname{Eq}(C)$ equipped with its canonical positive
determinant transport, as in the map on 15:48. It does not use every
possible rigidification in $\operatorname{Eq}(C,T)$.
Precomposition by the nontrivial Brandt class acts
by $f$. The quotient
$\mathcal T=\mathcal T'(C)/\langle f\rangle$ is a right
$H^1(A_{\rm int},\mu_2)$-torsor, independent of $C$ through the specified
comparison maps. Frobenius acts trivially on this quotient.

\begin{proof}
Given one equivalence, composing on the right with an autoequivalence
is free and transitive on equivalences, with the stabilizers accounted
for as natural transformations. The local calculation A.2a computes
the resulting groups of isomorphism classes. Any two global equivalences
between choices of $C$ differ by a Brandt autoequivalence, and that
difference is killed by the quotient. Their induced maps of quotient
torsors are therefore equal and satisfy the cocycle identity.
This is the precise sense of independence ``à isom unique près'' in
15:42: among maps induced by global comparisons. An abstract torsor
still has translation automorphisms; there is no uniqueness among
all equivariant maps.
On the prime-to-$p$ Tate line Frobenius is multiplication by $p$, at
$p$ it is the odd Brandt class, and at infinity it acts trivially.
Its translation is thus $f$, which becomes zero in the quotient.
The action factors through $\mathbb Z/2$, so the same conclusion
holds for $\widehat{\mathbb Z}$, not just its dense subgroup.
\end{proof}

\paragraph{A.2c (Proved: what the comparison gerbe is).}
Fix the crystalline determinant rule of A.1d. Over $A''=A'\times\mathbb R$
let $h$ be the $\mu_2$-gerbe of determinant-rigidified splittings of
$B_p$, and let $h_{\rm std}$ be the square-root gerbe of the Tate line
on $A'$ and the nonneutral standard $\mu_2$-gerbe on $\mathbb R$.
Put
\[
 \mathfrak G(\xi_p)|_{A''}=\operatorname{Eq}_{\mu_2}(h,h_{\rm std}),
 \qquad
 \mathfrak G(\xi_p)|_{\mathbb Z_p}=B\mu_2.
\]
This is a neutral fppf $\mu_2$-gerbe over $A_{\rm int}$. The set of its
sections up to isomorphism is a torsor under $H^1(A_{\rm int},\mu_2)$.
With the comparison construction of A.2b it is $\mathcal T$.
Its geometric section is the Tate-module realization with the Weil
pairing at every $\ell\ne p$, the unit object of $B\mu_2$ at $p$,
and the positive quaternionic realization at infinity. Its image in
the crystalline factor torsor is $\xi_p$, not a numerical unit without
the choice of an origin in that torsor.

\begin{proof}
Away from $p$ the algebra is Azumaya and the Tate realization gives
an equivalence of the two determinant gerbes. The real realization in
A.4a below does the same over $\mathbb R$. At $p$ a unit object is
part of the stated definition. These furnish a global object, hence
neutralize the comparison gerbe. For any object $x$ of a banded abelian
gerbe, the sheaf $\operatorname{Isom}(x,y)$ is a torsor under its band,
and $y\mapsto\operatorname{Isom}(x,y)$ is inverse to twisting $x$ by
a torsor. This proves the assertion about sections and their right
action. To identify it with $\mathcal T$, choose in each crystalline
comparison orbit the parity-zero representative, normalizing its
determinant class relative to $\xi_p$. The neutral $p$-gerbe then
contributes its unit-squareclass torsor. If we start with the opposite
parity representative, precomposition changes all local comparisons
by $f$, precisely the element killed in A.2b. This proves independence
of that choice and the stated identification of right torsors.
This neutral comparison gerbe is \emph{not} a neutralization of the
quaternion Brauer class over $\mathbb Q_p$. The latter class is
ramified, whereas the $p$-component here was expressly defined to be
neutral. Confusing these two gerbes would contradict Serre's argument.
\end{proof}

\paragraph{A.4a (Proved: quaternionic realization).}
Choose $E_0/k$ supersingular and an $\mathbb R$-algebra isomorphism
$\beta:\operatorname{End}^0(E_0)\otimes\mathbb R\to\mathbb H$.
For an elliptic curve $E$ in the supersingular category put
\[
 V_E=\operatorname{Hom}(E_0,E)\otimes_{\mathbb Z}\mathbb R,
 \qquad \langle x,y\rangle=\beta(x^\vee y).
\]
Transport the right action along $\beta$. Then $V_E$ is a right
quaternionic line with positive Hermitian form, and
$f:E\to E'$ acts by $x\mapsto fx$. It has adjoint $f^\vee$ and
norm multiplier $\deg f$. Any two such realizations are isomorphic
by a unitary natural transformation, and exactly two such
transformations exist. They differ by the central sign $-1$.

\begin{proof}
The rank-one module assertion is the covariant Deuring equivalence
(\citealp[Theorem 42.3.2 and Remark 42.3.3]{Voight2021}).
We have $x^\vee x=\deg x>0$ for nonzero rational $x$; after extension
to $\mathbb R$ this is the positive definite reduced-norm form.
Quaternion conjugation corresponds to duality, so the displayed
formula is Hermitian. The identities
$(fx)^\vee(fy)=x^\vee f^\vee f y=(\deg f)x^\vee y$ and
$\langle fx,z\rangle=\langle x,f^\vee z\rangle$ prove the assertions
about morphisms.
Between two choices of $E_0,\beta$, Skolem--Noether supplies an
intertwiner in the corresponding real Hom line. All intertwiners
differ by a nonzero real scalar. Normalizing its positive norm leaves
exactly the two choices differing by sign. Conversely naturality at
$E_0$ forces every natural automorphism to commute with its full
quaternion endomorphism algebra, so it is a real scalar; being unitary
forces $c^2=1$. This proves existence, exact ambiguity and independence
of choices in the correct groupoid sense.
The target has Hermitian lines as objects and \emph{all} quaternionic
linear maps as morphisms. Requiring every morphism to be an isometry
would exclude every isogeny of degree greater than one.
\end{proof}

\paragraph{A.3a (Proved: the two meanings of the real sign).}
There are two classes of unpolarized determinant comparison at infinity,
a torsor under $H^1(\mathbb R,\mu_2)=\{\pm1\}$.
Positivity selects one of these classes: a quaternionic linear
isomorphism has strictly positive reduced norm and cannot identify
a positive norm form with its negative. The chosen class still has
automorphism group $H^0(\mathbb R,\mu_2)=\{\pm1\}$.
Consequently positivity fixes the \emph{class} of the real component
of $|\xi|$, but does not select one natural isomorphism between its
realizations. The warning ``mais il reste l'ambiguïté des signes''
in 15:52 is correct in this latter sense. A claim of uniqueness up to
\emph{unique} isomorphism for A.4 would be false because $-1$ is a
nonidentity unitary natural automorphism.

\paragraph{A.7a (Proved under the Tate identification in R3).}
Let $h_{\rm mot}$ be the Artin-and-Tate-rigidified motivic gerbe of
15:7, with its asserted band $H=H^0\times\mu_2$. Its $\mu_2$
component has image $[B_p]\in\operatorname{Br}(\mathbb Q)[2]$, hence
local invariants
\[
 \operatorname{inv}_v(h_{\rm mot,\mu_2})=
 \begin{cases}\tfrac12,&v=p,\infty,\\0,&v\ne p,\infty.\end{cases}
\]
The supersingular category determines this component at every place.
It does not determine the $H^0$-component. The torsor $\mathcal T$
records comparisons in degree one between specified realizations
of this degree-two obstruction; the cohomology class alone does not
determine its pointed comparison data. This is the precise, qualified
version of Lead A2.

\begin{proof}
After rigidifying Artin realization, the geometric endomorphisms of
a supersingular elliptic curve are available; after rigidifying its
determinant, the residual scalar band on its two-dimensional splitting
module is $\mu_2$. Its image under $\mu_2\to\mathbb G_m$ is the
splitting gerbe of its endomorphism algebra, namely $B_p$.
Over $\mathbb Q$, Kummer identifies $H^2(\mathbb Q,\mu_2)$ with
$\operatorname{Br}(\mathbb Q)[2]$. Ramification gives the listed
invariants, and the global Brauer exact sequence determines the
class. In particular a real two-dimensional fibre functor would embed
$\mathbb H$ in $M_2(\mathbb R)$, which is impossible: a nonzero
homomorphism would be injective by simplicity and, by equality of
dimensions, an isomorphism, but one algebra is division and the other
is not. This rechecks the part of R3 actually used.
Twisting a comparison by any local $\mu_2$-torsor changes its
degree-one datum while preserving the underlying Brauer class.
Therefore no converse recovery of the pointed torsor from that
class is possible. The unbounded-denominator witnesses in Part D
separately distinguish the torus component.
\end{proof}

\subsection{Rigidified Morita equivalences}

\paragraph{A.5a (Proved under condition (C), including nonzero Picard groups).}
Let the unrigidified Brandt categories of two equivalent Deuring
categories carry the symmetric Picard structure obtained from
commuting two-sided fractional ideals. Let $N$ be their reduced-norm
functor to $\operatorname{Pic}(A')$. The rigidified autoequivalence
category is the homotopy fibre
\[
 \{(P,\theta):\ P\text{ an invertible bimodule},\quad
          \theta:N(P)\xrightarrow{\sim}A'\}.
\]
An equivalence of the determinant-rigidified categories induces a
canonical equivalence of these Picard categories, in the coherent
2-categorical sense, and identifies the left and right actions on
the category of comparison equivalences. No condition
$\operatorname{Pic}(A')=0$ is needed for this assertion.

\begin{proof}
Write an equivalence as an invertible $(O',O)$-bimodule $L$ with its
determinant compatibility. It sends $P$ to
$L\otimes_OP\otimes_OL^{-1}$, with the evident middle rings inserted.
The norm line of this bimodule is
$N(L)\otimes N(P)\otimes N(L)^{-1}$, canonically $N(P)$.
Hence the rigidification transports without choosing a basis of any
line. Composition and its associativity are preserved by cancellation.
Replacing $L$ by $L\otimes Q$ changes this transport by conjugation
with $Q$. The symmetric commutativity constraint gives
$Q\otimes P\otimes Q^{-1}\simeq P$; the hexagon and unit identities
give its compatibility with tensor product and with a third choice
of $L$. Applying the norm functor gives the same cancellation on
rigidifications. This constructs the coherent comparison rather
than claiming absence of all natural automorphisms.
For a comparison $L$ the actual interchange is
$(LPL^{-1})\otimes L\simeq L\otimes P$, which identifies the two
torsor actions.

For an unramified banded $\mu_n$-gerbe the same statement holds over
\emph{every} commutative base ring: its band-preserving
autoequivalences are canonically the Picard stack of $\mu_n$-torsors.
On a local neutralization this is twisting by a torsor; the description
descends because the band is abelian. Its classes fit into
\[
0\longrightarrow A'^*/A'^{*n}\longrightarrow H^1(A',\mu_n)
\longrightarrow\operatorname{Pic}(A')[n]\longrightarrow0.
\]
Thus omitting the Picard torsion would give a false generalization.
For example $A'=\mathbb Z[\sqrt{-5}]$ has the nonprincipal invertible
ideal $I=(2,1+\sqrt{-5})$ with $I^2=(2)$, producing a nontrivial
pair $(I,I^2\simeq A')$ invisible to unit squareclasses.
\end{proof}

\paragraph{A.5b (Proved: the hypothesis that must survive base change).}
Condition (C) holds for maximal orders over Dedekind bases: locally,
two-sided ideals are powers of the radical at a ramified place and
central ideals at an unramified place; hence their products commute,
and local equality implies global equality. This checks the R1 input
used in A.5a. It proves the requested rigidification for such admissible
base changes and for all Azumaya/gerbe base changes. The phrase
``catégorie de Deuring'' in 15:16 has a Dedekind maximal-order
qualification separated from the more general module-category
conditions. Under the latter conditions the universal quantifier over
arbitrary ramified, nonflat base changes is false, as the following
actual base change proves.

\paragraph{A.5c (Refuted: arbitrary ramified base change under the module-category definition).}
The universal statement at 15:36 is false if ``Deuring category'' retains
the module-category conditions (b),(c) of 15:16 after an arbitrary
commutative base change. This remains false when the category over the
new base is the actual base change of a global maximal-order category,
with the base-changed positive determinant object. Its category of
determinant-rigidified autoequivalences need not be a Picard category.
We give a counterexample over $R=\mathbb F_9$ in which even the group
of isomorphism classes is noncommutative.

\begin{proof}
We take the maximal order $O\subset B_3$ of A.1a and let $C$ be its
category of right modules locally free of rank one over
$\operatorname{Spec}\mathbb Z$. Put
\[
 R=\mathbb F_3[\alpha]/(\alpha^2+1),\qquad B=O\otimes_{\mathbb Z}R.
\]
Since $2$ is invertible at $3$, the localization of $O$ has generators
$i,j$ with $i^2=-1$, $j^2=-3$, and $ij=-ji$. Consequently $B$ has
orthogonal idempotents
\[
 e_0=(1+i/\alpha)/2,\qquad e_1=(1-i/\alpha)/2,
 \qquad x=e_0je_1,\quad y=e_1je_0.
\]
These form an $R$-basis. Their multiplication is specified by
$e_0+e_1=1$, $x=e_0xe_1$, $y=e_1ye_0$, and all products of two
elements of $Rx+Ry$ being zero. In coordinates the product and the
base-changed reduced norm are
\begin{align*}
 (a,b,c,d)(a',b',c',d')
   &=(aa',bb',ac'+cb',bd'+da'),\\
 N(ae_0+be_1+cx+dy)&=ab.
\end{align*}
The displayed polynomial is the reduction of the quaternion reduced
norm, not the determinant of the four-dimensional regular
representation (which would be its square).
The centre of $B$ is exactly $R$: commuting with $e_0$ forces $c=d=0$,
and commuting with $x$ then forces $a=b$. In particular no central
radical elements are being omitted.

This is the actual base change of $C$. Indeed, for finite projective
$O$-modules $P,Q$, base change identifies
\[
 \operatorname{Hom}_O(P,Q)\otimes R
 \simeq \operatorname{Hom}_B(P\otimes R,Q\otimes R).
\]
Every $P$ is free of rank one locally at $3$, so $P\otimes R$ is a
free rank-one $B$-module. Conversely $O\otimes R=B$ occurs, and over
the field $R$ every module locally free of rank one over
$\operatorname{Spec}R$ is free. Thus $C_R$ is equivalent to the
one-object $R$-linear category with endomorphism algebra $B$, and to
the full category of rank-one $B$-modules. This verifies conditions
(b),(c) of 15:16 without replacing $C_R$ by an unrelated category.

We specify its determinant rigidification carefully. For a free
rank-one right $B$-module $P$, its norm line is the contracted product
of the torsor of bases of $P$ with $R$ under $N:B^*\to R^*$.
Changing a basis by $b\in B^*$ acts on this line by $N(b)$.
This construction is the base change of the reduced-norm-line
construction for $O$. Choose the object $X=O$ before base change;
its identity basis has norm one and identifies its determinant object
with the positive reference object $T$ of A.6a. It therefore identifies
$\nu(X_R)$ with $T_R$, with no additional scalar choice.

On the one-object model an $R$-algebra automorphism $\rho$ of $B$
preserving the polynomial $N$ gives an admissible autoequivalence
$F_\rho$. Equivalently it gives the invertible bimodule with one
action twisted by $\rho$, free of rank one on both sides.
Since $N\circ\rho=N$, its induced functor on the determinant object
is the identity in the preceding coordinates. We rigidify it by
$u_\rho=\operatorname{id}_{T_R}$. These are specified rigidifications,
not merely an assertion that its transported norm line happens to
be trivial. Products have the same identity rigidification.
A natural isomorphism $F_\rho\to F_\kappa$ is a unit $b\in B^*$
satisfying $b\rho(a)=\kappa(a)b$ for all $a\in B$; its map on
the determinant line is multiplication by $N(b)$. It respects these
chosen rigidifications if and only if $N(b)=1$.

Set $t=1+\alpha$. Its square is $-\alpha$, its fourth power is $-1$,
and its order in $R^*$ is eight; hence $t$ is a nonsquare. Define
two $R$-algebra automorphisms by
\[
 S(e_0)=e_1,\quad S(e_1)=e_0,\quad S(x)=y,\quad S(y)=x,
 \qquad
 T_t(e_i)=e_i,\quad T_t(x)=tx,\quad T_t(y)=y.
\]
The multiplication table verifies these assertions directly, and
both preserve $N$. Write $\rho=S\circ T_t$, $\kappa=T_t\circ S$.
For the unit $d=te_0+e_1$ we have
\[
 \kappa=\operatorname{Int}(d)\circ\rho,\qquad N(d)=t.
\]
If $b\rho(a)=\kappa(a)b$ for every $a$, then $d^{-1}b$ commutes
with $\rho(B)=B$. The centre calculation gives $b=cd$ for
$c\in R^*$, and therefore
\[
 N(b)=c^2t\ne1.
\]
Thus the unrigidified products have an isomorphism, but no isomorphism
preserves their specified determinant rigidifications. The classes
of $(F_S,1)$ and $(F_{T_t},1)$ do not commute. A Picard category
necessarily has commutative $\pi_0$, so the proposed Picard structure
and hence the assertion of a canonical equivalence \emph{as Picard
categories} fail for this base change. Taking
$(C',T')=(C_R,T_R)$ already supplies the counterexample.

The script \texttt{verify\_ramified\_picard\_counterexample.py} checks
the multiplication, the automorphisms and the centre, enumerates all
$9^4=6561$ elements and all $8^2\cdot9^2=5184$ units, and finds
precisely the eight intertwiners $cd$. Their norms are the four
nonsquares of $R^*$; none is one. The argument above proves the result
independently of this finite verification.
\end{proof}

\paragraph{A.5d (Proved qualification of the counterexample).}
The category in A.5c satisfies the broad module-category definition
of 15:16. The algebra $B$ is not Azumaya: its nonzero radical
$Rx+Ry$ has square zero. It therefore does not satisfy the separately
bracketed maximal-order condition (a) over this field. If the phrase
``Deuring category'' is restricted to strict maximal-order categories
or to Azumaya categories after base change, A.5c is outside that
restricted domain and A.5a applies there. The universal quantifier
over arbitrary commutative base rings can thus be retained only
with an explicit admissibility condition on the resulting category;
under conditions (b),(c) alone it is refuted, rather than unresolved.


\subsection{Counting curves, counting the restricted triples}

\paragraph{A.8a (Known with derivation; literal identification refuted).}
Let $h(D)$ denote the class number of the imaginary quadratic order
of discriminant $D$. For $p>3$ the number $N_p$ of supersingular
elliptic curves over $\mathbb F_p$ up to $\mathbb F_p$-isomorphism is
\[
 N_p=\begin{cases}
 h(-4p),&p\equiv1\pmod4,\\
 4h(-p),&p\equiv3\pmod8,\\
 2h(-p),&p\equiv7\pmod8.
 \end{cases}
\]
The number of their distinct $j$-invariants is $N_p/2$.
For $p=2$, the two numbers are $3,1$; for $p=3$, they are $4,1$.
The numbers with trace zero are respectively $1$ and $2$ at these
small primes. Consequently neither counting $j$-invariants nor
counting all supersingular $\mathbb F_p$-forms counts the literal
maximal-$R$, trace-zero triples of 15:56.

\begin{proof}
Let $s$ be the number of geometric supersingular curves and $r$ the
number of their Frobenius-fixed $j$-invariants. The oriented Deuring
correspondence identifies maximal-order types with Frobenius orbits,
so the type number is $(s+r)/2$. The Deuring type-number formula
(\citealp[Proposition 30.9.2]{Voight2021}) says it is
$s/2+(h(-4p)+[h(-p)])/4$, where the bracket is zero if
$p\equiv1\pmod4$. Thus $r=(h(-4p)+[h(-p)])/2$.
Each such $j$ has two $\mathbb F_p$-forms. For a generic $j$ this is
quadratic twisting; for $j=0$ in the supersingular case the count is
$\gcd(6,p-1)=2$, and for $j=1728$ it is $\gcd(4,p-1)=2$.
The quadratic-order class-number relation is
$h(-4p)=3h(-p)$ for $p\equiv3\pmod8$, $p>3$, and
$h(-4p)=h(-p)$ for $p\equiv7\pmod8$ (ibid., 30.9.10).
These equalities give the result.

For $p=2,3$, the program enumerates all nonsingular general
Weierstrass equations and all changes
$x=u^2x'+r$, $y=u^3y'+su^2x'+t$. Every elliptic curve has such an
equation and every pointed isomorphism has this form. It partitions
the full finite set into orbits and verifies constancy of point counts
on each orbit. Its supersingular traces are $0,-2,2$ at $2$, and
$0,0,-3,3$ at $3$. The supersingularity criterion is $p\mid t$.
This is a finite exhaustive proof, not statistical evidence.
\end{proof}

\paragraph{A.8b (Proved neighboring unit statement; unfinished question unresolved).}
For $p>3$, $O^*=\{\pm1\}$ except at $j=0$, where its order is $6$,
and $j=1728$, where its order is $4$. These curves are supersingular
precisely for $p\equiv2\pmod3$ and $p\equiv3\pmod4$ respectively.
At $p=2$ the unique geometric supersingular curve has $24$
automorphisms; at $p=3$ it has $12$. The displayed ``respectively''
in the supplied transcription of 15:70 reverses the orders $4,6$.
These statements follow from the standard Weierstrass changes: for
odd characteristic greater than three an automorphism scales
$(x,y)$ by $(u^2,u^3)$, forcing $u^4=1$ when $j=1728$ and $u^6=1$
when $j=0$. The supersingularity assertions also follow from the
Hasse invariant coefficient or the inertness of $p$ in the quadratic
CM field. At $2$, use $y^2+y=x^3$: the changes have $u^3=1$,
$r=s^2$, $s^4=s$, and $t^2+t=s^6$, giving $3\cdot4\cdot2=24$.
At $3$, use $y^2=x^3-x$: one has $s=t=0$, $u^4=1$ and $r^3=r$,
giving $4\cdot3=12$.
The transcription ends at ``Question: $O^*$ est-il tjrs égal''.
Since its preceding diagram already states the generic unit answer,
we cannot identify the missing continuation with that same assertion.
Lead A3's guess about the missing words is not established by the
available witness. The original scan or a continuation is the exact
remaining evidence needed.
Michel Hua's project independently records that the ``page s'arrête
sur une question'' and repeats the same unfinished words in its
\texttt{15.modern.tex}. We credit that witness for confirmation of
the truncation only; it does not supply the missing assertion.

\paragraph{A.8c (Known with derivation: the corrected maximal-order count).}
Restrict to geometric trace-zero descents $E/\mathbb F_p$ with
$R\hookrightarrow\operatorname{End}_{\mathbb F_p}(E)$ sending
$\sqrt{-p}$ to Frobenius. Their number is
\[
 \#\operatorname{Pic}(R)=
 \begin{cases}h(-4p),&p\equiv1\pmod4,\\
 h(-p),&p\equiv3\pmod4,\\
 h(-8)=1,&p=2.
 \end{cases}
\]
In particular the answer is $h(-3)=1$ at $p=3$.
To derive this, trace zero is an elliptic isogeny class by
\citet[Theorem 4.1(5)]{Waterhouse1969}. Its rational endomorphism
algebra over $\mathbb F_p$ is $\mathbb Q(\sqrt{-p})$.
The maximal order is an allowed endomorphism ring by that paper's
Theorem 4.2(3). Theorem 4.5 then makes its isomorphism classes a
principal homogeneous space under $\operatorname{Pic}(R)$.
The embedding is fixed by its prescribed value on $\sqrt{-p}$, so
it adds no multiplicity. When $p\equiv3\pmod4$, the other
$h(-4p)$ trace-zero curves have the nonmaximal Frobenius order.
This also derives the small-prime trace-zero counts independently
of the exhaustive enumeration.

\paragraph{A.1g (Proved: a corrected equivalence at the level of rules).}
Although the claimed equivalence of \emph{all descent objects} is
false, restriction of the geometric determinant rules of A.1d to
the nonempty maximal-$R$ locus of A.8c is a bijection onto rules
compatible with its isomorphisms and invertible $R$-ideal action.
Indeed the ideal action is transitive by Waterhouse's Theorem 4.5,
so the ratio of two rules is a single unit squareclass. The
polarization rule restricts to such a rule, giving an origin on both
sides and identifying restriction with the identity on
$\mathbb Z_p^*/\mathbb Z_p^{*2}$. Its compatibility follows from
the positive norm of an isogeny, as in A.1d. The determinant
normalization in this description is really
$\bigwedge^2_{\mathbb Z}P_0\otimes
(\bigwedge^2_{\mathbb Z}R)^{-1}$, the quadratic norm line.
Writing it simply as $\det_{\mathbb Z}P_0\simeq\mathbb Z$ requires
the fixed ordered basis $(1,\sqrt{-p})$, or
$(1,(1+\sqrt{-p})/2)$ in the odd discriminant case, of $R$.
The norm-line formulation avoids an unmentioned orientation choice.
This statement covers $p=2,3$ as well; it does not identify their
additional nonzero-trace supersingular forms with these triples.

\paragraph{A-background audit.}
R1 is used only with condition (C) and coherent, not literal unique,
equivalences. R2 is checked in A.2a. R3 is conditional on the stated
Tate identification; its quaternion obstruction is checked in A.7a.
R4's norm assertion is checked in A.6a. Its separate classification
of all strict Deuring categories solely by $\operatorname{Br}(K)$
needs a degree qualification: over a field the rank-one module
categories for $M_2(K)$ and $M_3(K)$ each have one isomorphism class
of objects, with nonisomorphic endomorphism rings, so are not equivalent
despite having the same Brauer class. The fixed algebra $B_p$ used here
avoids that defect. None of these corrections is silently substituted
into the manuscript transcription.
